\section{Introduction}
%=============================================================
\subsection{Geometric Satake and numerical coincidences}

 In \cite{EQuantumI}, the geometric Satake equivalence \cite{LusztigGS, Ginz95, MirkVil} was reinterpreted algebraically. This reinterpretation makes sense in any type, but was only made explicit in type $A$, where we now focus our attention.

On one side of the equivalence is a 2-category built from the $\Q$-linear representation theory of $\mathfrak{sl}_n$, which we call $\Rep^{\Omega} \slf_n$. 
The objects of this $2$-category form a group $\Omega \cong \Z/n\Z$, which is isomorphic to the quotient of the weight lattice by the root lattice; the $1$-morphisms are representations whose weights live in the appropriate coset within $\Omega$, and $1$-morphism composition is tensor product of representations. Loosely speaking, $\Rep^{\Omega} \slf_n$ is just the monoidal category $\Rep \slf_n$ with additional bookkeeping.

On the other side of the equivalence is a 2-category of singular Soergel bimodules \cite{WillSingular} in type $\tilde{A}_{n-1}$, which can be thought of as an algebraic replacement for equivariant perverse sheaves on the affine Grassmannian. 
Crucial to note is that the vertices of the affine Dynkin diagram $\tilde{A}_{n-1}$ are in bijection with $\Omega$, and these vertices also parametrize copies of the finite Weyl group $S_n$ inside the affine Weyl group $W_{\aff}$ as (standard) parabolic subgroups. 
These finite parabolic subgroups are objects in the 2-category of singular Soergel bimodules $\SSBim$. The $1$-morphisms are particular bimodules over certain polynomial rings.

The geometric Satake equivalence becomes a (conjectural) $\Q$-linear 2-functor $\GS$ from $\Rep^{\Omega} \slf_n$ to $\SSBim$. 
Note that singular Soergel bimodules are a graded category whereas representations are not, so one cannot expect $\GS$ to be an equivalence. Instead it is a \emph{degree zero equivalence}: it is fully faithful onto the space of degree zero 2-morphisms, and there are no 2-morphisms of negative degree between the 1-morphisms in its image.
We call this the \emph{Soergelified (geometric) Satake equivalence}. We refer to it as conjectural for two reasons. First, the reformulation from traditional geometric Satake to Soergelified Satake was sketched but not made rigorous in \cite{EQuantumI}. Second, even were it made rigorous, geometric Satake would imply the existence of a $2$-functor $\GS$, but it would not be obvious how to make this $2$-functor explicit on the level of $2$-morphisms.

In \cite{EQuantumI} an explicit (potential) construction of the $2$-functor $\GS$ is given for all $n \ge 2$, and it is proven correct for $n=2, 3$. The construction uses webs, which are a generators-and-relations description of the monoidal category $\Rep \slf_n$. For $n=2$ webs form the well-known Temperley-Lieb category; for $n=3$ they were constructed by Kuperberg \cite{Kupe}; for $n>3$ they were constructed by Cautis-Kamnitzer-Morrison \cite{CKM} based on earlier work of Morrison. Concretely, \cite{EQuantumI} explains where the generating webs should go under $\GS$, but only checks the relations for $n=2,3$. In this paper we check the relations for $n > 3$.

Amongst the relations defining the web category, various signed binomial coefficients appear. For example, the value of a $1$-labeled circle in the web category is $(-1)^{n-1} \binom{n}{1} = (-1)^{n-1} n$, which is the categorical dimension of the standard representation of $\slf_n$. Meanwhile, the construction of singular Soergel bimodules takes as input the reflection representation $V$ of $W_{\aff}$, equipped with a choice of simple roots and coroots. 
The technical heart of Soergelified Satake is the realization that these binomial coefficients are also coded within $V$. The simplest example of this phenomenon is the fact that $-2$, the value of the Temperley-Lieb circle, agrees with the off-diagonal entry in the Cartan matrix for $\tilde{A}_1$.

The representation $V$ is spanned by the simple roots $\{\alpha_i\}_{i \in \Omega}$, and the action of $W_{\aff}$ is governed by the affine Cartan matrix. Let $R = R_{\aff}$ be the polynomial ring whose linear terms are $V$. For each simple reflection $s_i$, $i \in \Omega$, one defines Demazure operators $\partial_i \colon R \to R$,
\begin{equation} \partial_i(f) = \frac{f - s_i(f)}{\alpha_i}. \end{equation}
For example, the entries of the Cartan matrix are $\partial_i(\alpha_j)$; the number $-2$ from the previous paragraph is $\partial_1(\alpha_0)$ when $n=2$. More generally, one can realize binomial coefficients by applying certain sequences of Demazure operators to certain products of positive roots. For example, when $n=4$, one has
\begin{equation} \label{4intro} \partial_1 \partial_2 \partial_3 \left(\alpha_0 \cdot s_1(\alpha_0) \cdot s_2 s_1(\alpha_0) \right) = -4 = (-1)^{4-1} \binom{4}{1}. \end{equation}
The reader can guess the generalization of \eqref{4intro} to any $n \ge 2$. Similarly, 
\begin{equation} \label{42intro} \partial_2 \partial_3 \partial_1 \partial_2(\alpha_0 \cdot s_1(\alpha_0) \cdot s_3(\alpha_0) \cdot s_1 s_3(\alpha_0)) = \binom{4}{2}.
\end{equation}
To our knowledge, these observations and related numerical coincidences have not yet appeared in the literature.

The computations \eqref{4intro} and \eqref{42intro} arise when checking that the circle relations of the web category continue to hold in $\SSBim$ after applying $\GS$. The product of roots appears via a coproduct for a Frobenius extension between subrings of $R$, and the sequence of Demazure operators is the trace map for a Frobenius extension between different subrings of $R$. To check the remaining web relations requires more computations of a similar flavor: applying Demazure operators to tensors (of polynomials) built from Frobenius coproducts. These computations can be quite thorny.

Once one has constructed the functor $\GS$, one can prove the rest of the Soergelified Satake equivalence from scratch (i.e. without going through geometric Satake). It remains to show that $\GS$ is a degree zero equivalence, and here is a brief sketch. Williamson's Hom formula (\cref{thm-hom formula}) shows that the graded dimension of $2$-morphism spaces between singular Soergel bimodules is determined by a certain pairing on its Grothendieck group. The seminal work of Lusztig \cite{LusztigGS} determines the images of (characters of) irreducible representations under the Satake isomorphism (i.e. the decategorification of the geometric Satake equivalence).  These results combine to show that the dimensions of the $2$-morphism spaces are appropriate for $\GS$ to be a degree zero equivalence. In particular, if $\GS$ is faithful then it is also full (to degree zero) by a dimension count.

If we define our categories to be $\Q$-linear, then we can continue as follows. The semisimple 2-category $\Rep^{\Omega} \slf_n$ has no nontrivial monoidal ideals (see \cref{lemma monoidal ideals in qRep} for a proof, but this is an adaptation of a well-known argument), so the kernel of $\GS$ is either zero or everything. The latter is absurd, since identity maps go to identity maps. Hence $\GS$ is faithful. This proof will not work integrally or over fields of finite characteristic, where nontrivial monoidal ideals exist; we will have more to say about this problem later.

% \footnote{Let us spell out this well-known fact; the reader can ignore $\Omega$ and just work with the more familiar monoidal category $\Rep \slf_n$. Because $\Rep \slf_n$ is semisimple, any 2-sided ideal is generated by the identity maps of a set of simple objects. However, since the trivial representation $L_0$ is a summand of $V \ot V^*$ for any representation $V$, a monoidal ideal containing $\id_V$ will also contain $\id_{L_0}$. Conversely, a monoidal ideal containing $\id_{L_0}$ will contain $\id_{L_0 \ot V} = \id_V$ for all representations $V$.}

\subsection{The quantum deformation}

One side of the Satake equivalence has a well-known $q$-deformation: the representation theory of the quantum group $U_q(\slf_n)$. These can also be described with webs, and quantum binomial coefficients replace ordinary binomial coefficients. Interestingly, \cite{EQuantumI} introduced a $q$-deformed Cartan matrix of type $\tilde{A}_{n-1}$, leading to a deformations $V_q$ and $R_q$, and to a deformation of singular Soergel bimodules which currently has no analogue in geometry. The \emph{quantum (Soergelified) Satake equivalence} is the corresponding (conjectural) $\C(q)$-linear degree zero equivalence $\GS_q$ from $\Rep^{\Omega} U_q(\slf_n)$ to $q$-deformed singular Soergel bimodules. Here is the deformed Cartan matrix in type $\tilde{A}_3$.
\begin{equation} \label{eq:deformedq} \begin{bmatrix}
    2 &-1 &0 & -q^{-1} \\
    -1 &2 &-1 &0\\
    0 &-1 &2 &-q \\
    -q & 0&-q^{-1} &2
\end{bmatrix} \end{equation}
Repeating the computation of \eqref{4intro} for this Cartan matrix, one obtains
\begin{equation} \label{q4intro} \partial_1 \partial_2 \partial_3 \left(\alpha_0 \cdot s_1(\alpha_0) \cdot s_2 s_1(\alpha_0) \right) = -[4]_q. \end{equation}

Once again, \cite{EQuantumI} explicitly constructed $\GS_q$ for $n=2,3$ and proved it was well-defined. The construction of $\GS_q$ for $n \ge 4$ as a deformation of $\GS$ entails a choice of invertible scalar for each generating web, such that the web relations hold after applying $\GS_q$. In this paper we provide those scalars and check the relations, thus constructing $\GS_q$ for $n \ge 4$. We also prove that this $2$-functor is a degree zero equivalence over $\Q(q)$, finishing the work begun in \cite{EQuantumI}. The details can be found in \cref{thm main}. More accurately, we work over an extension of $\Q(q)$, see \S\ref{subsec:covintro}.

The technical heart of our proof is a number of thorny computations with Demazure operators acting on $R_q$. This requires even more care than one might expect, because the $q$-deformed Demazure operators no longer satisfy the braid relations, so one must keep careful track of Frobenius structures and transformations between reduced expressions in addition to the other algebraic work.

Our proof that $\GS_q$ is a degree zero equivalence over $\Q(q)$ follows the same outline as for $\GS$ in the previous section. However, Williamson did not prove his Hom formula in this level of generality. The representation $V_q$ is not \emph{balanced} (see \S\ref{subsec-balanced}), a technical condition without which one cannot pick out positive versus negative roots consistently. The Hom formula is one of several basic facts about Soergel bimodules which had not been proven for the $q$-deformation, a situation we rectify in this paper. This also allows us to use certain shortcuts: knowing the dimensions of morphism spaces allows us to bypass certain computations. 

Since our goal is to establish the quantum geometric Satake equivalence rigorously, and since the literature has a tendency to wave its hands at some technicalities, we will be concrete and fill in many details. So that later work can more easily establish the integral or finite characteristic versions of (quantum) Soergel Satake without redoing fundamental work, we prove our results over general base rings when possible. For example, we introduce reflection faithfulness (another technical condition) over general base rings, and show that the $q$-deformed reflection representation is reflection faithful. Along similar lines, the literature on geometric Satake tends to use the extended affine Weyl group, whereas we prefer using different parabolic subgroups in the affine Weyl group. There are some frustrations in matching the choices of convention (coming from right actions vs. left actions\footnote{The commutativity of the spherical Hecke algebra has obscured the effects of some conventional issues.}, positive vs. negative Iwahori subgroups, etcetera), so we redevelop some of the combinatorics from scratch. We provide a description of the combinatorial correspondence between dominant weights and double cosets in a fully explicit fashion, focusing on particular representatives of these cosets which are best adapted to the concept of a ``singular reduced expression'' as introduced in \cite{WillThesis, EKo}.

Another goal of this paper was originally to rigorously establish the relationship between Soergel Satake and geometric Satake (in the non-quantum setting). This is work in progress and should hopefully appear in an appendix to a future version of this paper. The relationship is sketched in \cite[\S 6]{EQuantumI} but no attempt is made at a rigorous proof there.  The reader who is curious about the relationship is encouraged to read the introduction of \cite{EQuantumI}, which instructs the geometrically-minded reader on how to skip most of the paper and efficiently reach \S 6.


\subsection{Remarks on integral forms}\label{subsec:youthfuloptimism}

In \cite{EQuantumI} it was promised that the general case of quantum Soergel Satake would be proven in a sequel. The twelve year time lapse between \cite{EQuantumI} and its sequel (the present paper) deserves some explanation. In fact, what \cite{EQuantumI} proposed is a stronger result about the integral forms of these categories, which we discuss now. We do not prove this stronger result.

\begin{remark} The first author apologizes for their youthful optimism; the intervening years have cured both these flaws, youthfulness and optimism. At the time of \cite{EQuantumI} the basic computations like \eqref{q4intro} had been done (and \eqref
{42intro} for $n=4$), but not the nastier ones. \end{remark}

By presenting a category by generators and relations, one makes it straightforward to define functors therefrom. Another aspect of a presentation is that it defines an \emph{integral form} of the category. The category one is presenting may originally be $\Q$-linear (or $\Q(q)$-linear), but if all the coefficients in all the relations live in a subring $A$ (like $\Z$ or $\Z[q,q^{-1}]$), then one can define an $A$-linear category by generators and relations instead. We think of an integral form as being ``correct'' if morphism spaces are free over $A$, and specialize in an interesting way (e.g. to finite characteristic, or setting $q$ to be a root of unity). It was proven in \cite{ELLCC} that Cautis-Kamnitzer-Morrison's webs are the correct integral form for $\Rep^{\Omega}\slf_n$, in that they specialize to describe the category of tilting modules over any field.

Meanwhile, it was already known by the authors of \cite{CKM} that many of their relations were redundant when working non-integrally, over $\Q(q)$. For example, a streamlined presentation can be found in work of Bigelow \cite{Bigelow} (which also fixes a sign error in \cite{CKM}), and we reiterate that Bigelow's presentation is knowingly the ``incorrect'' integral form.

Our construction of $\GS$ and $\GS_q$ factors through a diagrammatic description of $\SSBim$ which we denote $\DiagSBS$. While an integral form for ordinary Soergel bimodules was found in the early 2010s by Elias, Khovanov, and Williamson \cite{EKho, ECathedral, EWGr4sb}, constructing an integral form for singular Soergel bimodules remains a long-standing programme (one which the youthful optimist expected to have been concluded by now). 
A diagrammatic $2$-category was constructed in \cite{EWSFrob}, called $\DiagSBS(\Gamma)$ in the body of the paper, with a 2-functor to $\SSBim$ which is full (proven\footnote{in the case when the realization is balanced.}  in \cite{EKLP}) but known \emph{not} to be faithful; there are not enough relations in \cite{EWSFrob}. Let $\DiagSBS$ be the quotient of $\DiagSBS(\Gamma)$ by the kernel, so that $\DiagSBS$ maps fully faithfully to $\SSBim$. Some of the relations in this kernel are known from the literature (e.g. from \cite{ECathedral, ELosev}), but it is expected that more relations are required, especially to get the ``correct'' integral form. Thus $\DiagSBS$ is not very explicit in its current state, but is nonetheless a useful language to encode morphisms in $\SSBim$.

One of the ``shortcomings'' of this paper is that we do not prove our 2-functor $\GS_q$ is well-defined using only known relations in $\DiagSBS$, but resort to calculations in $\SSBim$. We cannot be certain, but feel there are not enough known relations to do the job, so we await an improvement in the available diagrammatic technology. Since the state of the art constrains us to work with $\SSBim$ and not to work with an integral form thereof, we avail ourselves of certain tricks. We use the Soergel-Williamson Hom formula (which specifies the dimension of certain morphism spaces) to simplify some computations. We avail ourselves of Bigelow's presentation.

We have been referring to the (not-explicitly-defined) $2$-category $\DiagSBS$ as an integral form for $\SSBim$, but this is abuse of terminology. It should be an integral form for the \emph{Hecke $2$-category}, which is supposed to categorify the Hecke algebroid in such a way that the Soergel-Williamson Hom formula holds. It is known that singular Soergel bimodules are \emph{not} well-behaved in some settings (e.g. finite characteristic in affine type) so they are only a good model for the Hecke $2$-category generically, not for its integral form. Abe \cite{Abe} has provided an algebraic model for the Hecke $2$-category, though not a presentation.

When the programme to present the Hecke $2$-category is complete, there should be a version of $\DiagSBS$ defined over $\Z[q,q^{-1}]$, with a presentation, such that morphism spaces are \textbf{free} over $\Z[q,q^{-1}]$ of the expected size. Showing freeness given a presentation is very difficult work! What it would take to upgrade our $2$-functor $\GS_q$ to a $2$-functor from Cautis-Kamnitzer-Morrison's presentation to the integral form of $\DiagSBS$? What remains is less difficult, though there are some subtleties.

We have defined the $2$-functor on objects and morphisms already, and need only check the relations. Most of the relations from \cite{CKM} have the form: take a diagram $X$, and rewrite it as a linear combination of elements in a basis $\mathbb{B}$ for this morphism space. Consider the collection of diagrams $\GS_q(\mathbb{B})$ in $\DiagSBS$. This collection is linearly independent, which can be proven after base change to $\Q(q)$ where it follows from our results here. If $\GS(X)$ could be rewritten as a linear combination of $\GS_q(\mathbb{B})$, then the coefficients must be as desired, again provable after base change to $\Q(q)$. However, it is not obvious that $\GS_q(\mathbb{B})$ is a basis, even though it has the desired size; it might only span a $\Z[q,q^{-1}]$-sublattice. What remains is not a computational problem, but an abstract question: is $\GS_q$ full?

Major advances in the aforementioned programme were recently made by the first author together with Ko, Libedinsky, and Patimo \cite{EKo, KELPDemazure, KELPBruhat, KELPLeibniz, EKLP}. In particular, in \cite{EKLP} a basis for 2-morphisms in $\SSBim$ is constructed combinatorially as the image of certain diagrams in $\DiagSBS(\Gamma)$, called the \emph{double leaves basis}\footnote{It is proven in \cite{EKLP} that the double leaves maps become a basis after applying the 2-functor to $\SSBim$, but this is done under the balanced assumption. More work would be needed to adapt their arguments to the $q$-deformation.}. This is expected to be a basis for the correct integral form of $\DiagSBS$, once that integral form is established. Meanwhile, a combinatorial basis for webs was constructed in \cite{ELLCC}, the \emph{double ladders basis}. One way to prove the fullness of $\GS$ is to prove the following.

\begin{conjecture} \label{conj:double} Under $\GS$ (or $\GS_q$), double ladders are sent to double leaves (up to invertible scalar in $\Z[q,q^{-1}]$). \end{conjecture}

Double ladders are built from elementary light ladders, which live in one-dimensional morphism spaces modulo lower terms. That the image of elementary light ladders agrees with a (non-elementary) light leaf modulo lower terms is thus not surprising, but in the examples we have computed they agree on the nose.

At first glance this conjecture seems straightforward, but it is surprisingly complex. We illustrate with an example, and though we have not yet introduced the diagrammatic categories in question, we expect the reader will get the point. Complete details on this example are found in \cref{subsec:LLappendix}, so pardon the incompleteness here.

A relatively simple form of elementary light ladder is exemplified by the following web. We also show where it goes under $\GS$.
\begin{equation} \label{nastyELL}\isvg{0.15}{ELLforBen2} \quad\mapstosize{1.5} \quad {
\labellist
\small\hair 2pt
 \pinlabel {$\blk{5}$} [ ] at 8 -5
 \pinlabel {$\blk{6}$} [ ] at 24 -5
 \pinlabel {$\blk{2}$} [ ] at 40 -5
 \pinlabel {$\blk{5}$} [ ] at 56 -5
 \pinlabel {$\blk{0}$} [ ] at 72 -5
 \pinlabel {$\blk{2}$} [ ] at 88 -5
 \pinlabel {$\blk{4}$} [ ] at 8 80
 \pinlabel {$\blk{6}$} [ ] at 24 80
 \pinlabel {$\blk{0}$} [ ] at 72 80
 \pinlabel {$\blk{4}$} [ ] at 88 80
 \pinlabel {$\dgr{0}$} [ ] at 90 45
 \pinlabel {$\dgr{6}$} [ ] at 2 45
\endlabellist
\centering
\ig{1}{NastyELLimage}
} \end{equation}
\vspace{.2cm}
Meanwhile, the corresponding light leaf in $\DiagSBS$ is the following diagram.
\begin{equation} \label{nastyLL} \ig{1}{NastyLL} \end{equation}
It turns out that \eqref{nastyELL} and \eqref{nastyLL} are actually equal, though this requires a number of non-obvious diagrammatic transformations (e.g. ``switchback relations'').

\subsection{A change of variables} \label{subsec:covintro}

We wish to note a change of variables that is used throughout the paper. The $q$-deformed Cartan matrix for $\tilde{A}_n$ from \eqref{eq:deformedq} has the nice feature that it contains the usual Cartan matrix for $A_n$ inside, but it admits very little symmetry! Each parabolic subgroup of $\tilde{A}_n$ behaves a little bit differently. In the body of this paper, we do most computations for a different deformation of the usual $\tilde{A}_n$ Cartan matrix over a different formal variable $\zeta$, exemplified here for $n=3$.
\begin{equation} \begin{bmatrix}
    2 &-\zeta &0 &-\zeta^{-1} \\
    -\zeta^{-1} &2 &-\zeta &0 \\
    0 &-\zeta^{-1} &2 &-\zeta \\
    -\zeta & 0 &-\zeta^{-1} &2 
\end{bmatrix}\end{equation}
This Cartan matrix was introduced in \cite{EJY1}. It leads to deformations $V_{\zeta}$ and $R_{\zeta}$. We will also define a $2$-functor $\GS_{\zeta}$ from $\Rep^{\Omega} \slf_n$ to $\SSBim$ as defined over $R_{\zeta}$. 

When $\zeta^n=q^{-2}$, there is an isomorphism $V_q \cong V_{\zeta}$ which intertwines the root data up to scalar. This change of variables requires an extension of $\Q(q)$. The extension $\Q(q^{\frac{1}{n}})$ will suffice, and we use this for the introduction, see \cref{subsubsec quantum numbers} for more details.

The advantage of the $\zeta$-deformation is its obvious rotational symmetry (i.e. rotation of the Dynkin diagram, the action of $\Omega$; not rotation of diagrams), a symmetry we denote $\sigma$ in this paper, and a symmetry shared by colored webs. Diagrammatically, $\sigma$ will just change various labels (depicted as colors) in each diagram, without changing the underlying diagrams or rescaling them. The 2-functor $\GS_{\zeta}$ will preserve $\sigma$, making it easier to check relations.

\begin{remark} One could translate the symmetry $\sigma$ through the isomorphism with $V_q$ to obtain a symmetry on $V_q$ and on categories built from $V_q$. However, this symmetry would not just relabel things but would also multiply them by annoying scalars. \end{remark}

The disadvantage of the $\zeta$ deformation is that it obscures the quantum parameter $q$ which is most common in the literature. Knowing the isomorphism $V_q \cong V_{\zeta}$ is not quite enough to easily convert from one parameter to another. The most natural choice of Frobenius extension structures between $R_q$ and its subrings (all of which can be defined over $\Q(q)$ without needing an extension) differs from the most natural choice for $R_\zeta$. The very meaning of the diagrams we use (in $\DiagSBS$) depends on this choice of Frobenius extension structure. We distinguish the two diagrammatic categories as $\DiagSBS_\zeta$ and $\DiagSBS_q$. In \cref{appendix:comparisonqz} we explicitly describe the way in which one converts between these two conventions.

One might also hope to define $\GS_q$ over $\Q(q)$ rather than its extension $\Q(q^{\frac{1}{n}})$, and this is possible, but we do not attempt to make it explicit, because of the following obnoxious issue.

As noted earlier, diagrams in the web category are sent to scalar multiples of diagrams in $\DiagSBS$. We offer a particular preferred choice of these scalars in \cref{goodchoiceoflambdanu}. What precisely these scalars are is not intrinsic to geometric Satake! One could obtain a new $2$-functor $\GS'$ with different scalars by precomposing with an object-fixing autoequivalence of webs, one which rescales the generating morphisms. However, certain relationships between these scalars are required for the $2$-functor to be well-defined, see \eqref{scalarconditions}. When we define $\GS_{\zeta}$ in \cref{defn dQGS functor}, we assume implicitly that these scalars will be preserved by $\sigma$. Abandoning this assumption will lead to more bookkeeping.

If we naively try to define $\GS_q$ by using our preferred choice of scalars and pass through the equivalence between $\DiagSBS_{\zeta}$ and $\DiagSBS_q$, the resulting scalars will live in $\Q(q^{\frac{1}{n}})$ and not the subring $\Q(q)$. One must choose scalars which are not preserved by $\sigma$ in order to obtain a 2-functor defined over $\Q(q)$ (one must treat each parabolic subgroup a little differently). This is what was done for $n=3$ in \cite{EQuantumI}. Currently, the extra bookkeeping does not seem to be worth the effort, though we may spell this out in a future version of this paper.

\subsection{More symmetries} 


Web categories have three ``natural'' symmetries we focus on: rotation of diagrams by 180 degrees, the horizontal flip, and the vertical flip. The $2$-category $\DiagSBS(\Gamma)$ from \cite{EWSFrob} also has these symmetries\footnote{In $\DiagSBS(\Gamma)$, each strand in a diagram has an orientation, but this orientation is cosmetic, a convention to help one keep track of a labeling on regions. After flipping a diagram, convention dictates that one will also need to invert orientations.}. Note that the composition of a horizontal and a vertical flip is the rotation by 180 degrees. The 2-functor $\GS_q$ appears to preserve these symmetries; a flipped web is sent to a flipped diagram (up to sign). However, as noted in \S\ref{subsec:youthfuloptimism}, we are only able to construct a $2$-functor to singular Soergel bimodules $\SSBim$; not to $\DiagSBS(\Gamma)$, nor to an explicit quotient thereof which a priori possesses the same symmetries.

The bimodule $2$-category $\SBSBim$ also has three symmetries in the literature: adjunction, swapping the left and right actions, and (Poincar\'{e}) duality (in the same order as the corresponding diagrammatic symmetries above). Duality is the most complex of these, an exposition for which can be found in \cite[Section 3.4]{BBE}. For an $(R^I, R^J)$-bimodule $B$, its right dual (ignoring grading shifts) would be $\Hom_{(-,R^I)}(B,R^I)$, where this represents the space of morphisms as a right $R^I$-module. Any singular Bott-Samelson bimodule is isomorphic to its right (or left) dual, and that isomorphism is defined by fixing a particular invariant bilinear form on the underlying bimodule. Crucially, these invariant bilinear forms can be defined using same Frobenius extension structures which are used to define adjunction. It is for this reason that ``rotation'' agrees with the composition of the ``horizontal'' and the ``vertical'' flip in the category of bimodules. This connection between bilinear forms and Frobenius structures is sometimes explicit and sometimes implicit in the literature, but what is seemingly absent from the literature is the statement that the 2-functor $\Phi$ from diagrams $\DiagSBS(\Gamma)$ to bimodules intertwines these three symmetries, especially in the unbalanced case. We clarify the situation in \cref{subsec reversal duality and phi}, and state the compatibility between $\GS_q$ and these symmetries in \cref{lemma symmetry of GS-zeta}. 

% Meanwhile, adjunction is defined using Frobenius extension structures. From these descriptions it should not be obvious that the composition of left-right swap and duality agrees with adjunction (up to coherent isomorphism). However, this is indeed the case, because the bilinear form in question is also related to the Frobenius extension structure, though this connection is at best implicit in the literature, see \cref{rmk intersection forms in literature}. Relatedly, the literature also does not seem to contain the statement that the functor $\Phi$ from diagrams $\DiagSBS(\Gamma)$ to bimodules intertwines the three symmetries discussed above, especially not for the unbalanced case. With this statement in place, one
% concludes immediately that $\GS_q$ preserves these three symmetries \PT{Ummm... $\GS_\zeta$ doesn't commute with horizontal swapping..., though $\GS_q$ might, after we find the $q$-form, likely breaking color symmetry.}.

% These things we all make explicit in \cref{subsec reversal duality and phi}, though we leave many calculations to the reader. Once duality is cleanly stated, these indeed become fairly straightforward exercises. 


\subsection{Acknowledgments}

Both authors were supported by the first author's NSF grant DMS-2201387. We appreciate the support given to our research group by NSF RTG grant DMS-2039316. This work forms a substantial portion of the second author's doctoral thesis. The second author is extremely grateful to the first author, who as a doctoral advisor was a source of invaluable guidance, patience, and support. The second author also thanks Victor Ostrik for priceless mathematical discussions and encouragement throughout his doctoral studies.





\comm{

\section{Extended introduction}

\subsection{Reformulating geometric Satake}

The goal of \cite{EQuantumI} and the current paper is to directly prove the Soergelified Satake equivalence using straightforward algebra, rather than deducing it from the geometric Satake equivalence. After all, this is the only route currently available to prove the quantum version of the result. It is not the goal of either paper to provide a rigorous proof of the translation from geometric Satake to Soergelified Satake (and the equivalence between these two results). An outline of how this translation would work was given in \cite[Section 6]{EQuantumI}, though it relies on numerous folklorish details. For the sake of the geometrically-minded reader, we briefly recall the outline here. \BE{added:} We plan to add an appendix to this paper which contains more rigorous proofs of this geometric reformulation.

After the definition of $\Rep^{\Omega} \gf^{\vee}$, the rest of the section may be skipped harmlessly.

Geometric Satake is traditionally stated in the context of a semisimple algebraic group $G$ and its Langlands dual group $G^{\vee}$.
Let $\OC = \CC[[t]]$ and $\KC = \CC((t))$. Geometric Satake is an equivalence of monoidal categories
\begin{equation} \Rep G^{\vee} \to \Perv_{G(\OC) \times G(\OC)}(G(\KC)), \end{equation}
where the right-hand side indicates the category of $G(\OC) \times G(\OC)$- equivariant perverse sheaves with a convolution monoidal structure. The two actions of $G(\OC)$ come from left and right multiplication.
Meanwhile, Soergelified Satake is a relationship between structures attached to a complex semisimple Lie algebra $\gf$ and its Langlands dual $\gf^{\vee}$.
Depending on the choice of $G$ lifting $\gf$, there are several different geometric Satake equivalences that are encoded in the same Soergelfied Satake equivalence. 
We begin by assuming that $G^{\vee}$ is simply-connected, so that $\Rep G^{\vee} \cong \Rep \gf^{\vee}$.
In this case, $G$ has adjoint type. We discuss the general case in Remark \ref{rmk:notadjoint}.

Let $\Omega$ denote the weight lattice of $\gf^{\vee}$ modulo its root lattice, a finite abelian group. The category $\Rep \gf^{\vee}$ splits into blocks parametrized by $\Omega$, as each irreducible representation has weights living within a single coset for the root lattice. For $\chi \in \Omega$ we denote the corresponding block as $\Rep_{\chi}$. One can think of this decomposition as coming from central characters of $G^{\vee}$, as $\Omega \cong \Hom(Z(G^{\vee}), \C^\times)$. The block decomposition of $\Rep \gf^{\vee}$ means that any additive endofunctor of the category can be viewed as a matrix of functors, where each matrix entry would be a functor $\Rep_{\chi} \to \Rep_{\chi'}$ for some $\chi, \chi' \in \Omega$. For example, let $L$ be an irreducible representation in $\Rep_{\chi}$. The functor $L \otimes (-)$ is not irreducible, but splits as a direct sum of its nonzero matrix entries, one for each choice of $\nu \in \Omega$: this matrix entry is a functor which takes $\Rep_{\nu}$ to $\Rep_{\nu + \chi}$ via $L \otimes (-)$, and kills all irreducible representations not in $\Rep_{\nu}$.

A transparent way to encode this block decomposition is to replace $\Rep \gf^{\vee}$ with the 2-category $\Rep^\Omega \gf^{\vee}$. The objects of this $2$-category are parametrized by $\Omega$, and for each $\chi, \nu \in \Omega$, the category $\Rep_{\nu}$ is the 1-morphism category from $\chi$ to $\chi+\nu$. Composition of 1-morphisms is tensor product of representations. This is a full sub-$2$-category of $\Cat$, the $2$-category of categories, where each $1$-morphism $L \in \Hom(\chi,\chi+\nu)$ is thought of as the functor $L \otimes (-)$ between the appropriate blocks.

We return to geometry, where $G$ has adjoint type. We also have
\begin{equation} \Omega \cong \pi_1(G) \cong \pi_0(G(\KC)/G(\OC)) = \pi_0(G(\KC)). \end{equation}
For each $\chi \in \Omega$ we choose an element of $G(\KC)$ in the appropriate connected component, which we abusively call $\chi$. Left multiplication by $\chi$ is a map $\ell_{\chi} \colon G(\KC) \to G(\KC)$ which sends the nulcomponent isomorphically to the component of $\chi$.  If $\PC$ is a $G(\OC) \times G(\OC)$-equivariant perverse sheaf supported on the component of $\chi$, then $\ell_\chi^* \PC$ is supported instead on the nulcomponent, and is equivariant under $\chi^{-1} G(\OC) \chi \times G(\OC)$; equivariance for the left action has been twisted. 
Instead of allowing $G(\OC) \times G(\OC)$-equivariant sheaves on all components, we can consider sheaves only on the nulcomponent but can allow the equivariant group to change. After this change of perspective, to make convolution easier, it is more natural to also change the group acting on the right. That is, we consider the 2-category whose objects are $\Omega$, thought of as parametrizing ``conjugates'' of $G(\OC)$, and whose morphism category from $\chi_2$ to $\chi_1$ consists of the $(\chi_1^{-1} G(\OC) \chi_1, \chi_2^{-1} G(\OC) \chi_2)$-equivariant perverse sheaves on the nulcomponent of $G(\KC)$.

\begin{remark} \label{rmk:notadjoint} The nulcomponent of the affine Grassmannian is ``independent'' of which $G$ one chooses. That is, this nulcomponent agrees with the affine Grassmannian associated to the simply-connected group lifting $\gf$, see \cite[Theorem 1.3.11 (3)]{Zhu}. When $G$ is not of adjoint type, it still has $\Omega$ many distinguished ``conjugates'' of $G(\OC)$, though they may be conjugate under an outer automorphism rather than an inner one. When $G = SL_2$, these conjugates are
\begin{equation} P_0 := SL_2(\OC) = \begin{bmatrix} \OC & \OC \\ \OC & \OC \end{bmatrix}, \quad P_1 := \begin{bmatrix} t^{-1} & 0 \\ 0 & 1 \end{bmatrix} SL_2(\OC) \begin{bmatrix} t & 0 \\ 0 & 1 \end{bmatrix} = \begin{bmatrix} \OC & t^{-1}\OC \\ t \OC & \OC \end{bmatrix}. \end{equation}
Note that the diagonal matrix with entries $(t,1)$ is not an element of $SL_2(\KC)$, but it is an element of the adjoint type group $\PGL_2(\KC)$. Thus $P_1$-equivariant perverse sheaves on the nulcomponent of the affine Grassmannian agrees with perverse sheaves on a different component for $\PGL_2$, but is a category typically ignored when discussing geometric Satake for $SL_2$. Changing the equivariance group allows the affine Grassmanian to access all representations of $\gf^{\vee}$, regardless of the choice of $G$.
\end{remark}

The next step is to replace loop groups $G(\KC)$ with Kac-Moody groups $G_{\KM}$ in affine type. This replaces $G(\OC)$ with the parabolic subgroup $P_{\fin}$ associated to the finite Dynkin diagram inside the affine one. A vertex in the affine Dynkin diagram is \emph{removable} if removing it yields a diagram isomorphic to the finite Dynkin diagram. An important point is that there is also a natural bijection between $\Omega$ and the set of removable vertices in the affine Dynkin diagram \BE{of affine $\gf$ or of affine $\gf^{\vee}$ or of both...} (sending $0 \in \Omega$ to the affine vertex\footnote{Any vertex of the affine Dynkin diagram other than the affine vertex can be viewed as a vertex in the finite Dynkin diagram, and has an associated fundamental weight. The fundamental weights associated to removable vertices, together with the zero weight, enumerate the cosets in $\Lambda_{\wt}/\Lambda_{\rt}$. The first author thanks Pavel Etingof for this elaboration of the bijection.}). One replaces $\chi^{-1} G(\OC) \chi$ with a parabolic subgroup $P_{\chi}$ associated to a different copy of the finite Dynkin diagram, associated to another removable vertex. Replacing $G(\KC)$ with $G_{\KM}$ or $G(\OC)$ with $P_{\fin}$ will enlarge the size of the torus by an extra factor of $(\C^{\times})^2$, factors which cancel when considering the affine Grassmannian:  $G(\KC) / G(\OC) \cong G_{\KM}/P_{\fin}$. However, the larger torus in the equivariance group leads to slightly more data when considering perverse sheaves in the Kac-Moody setting, so this can be viewed as an upgrade. \BE{We should talk about your comment here}
\PT{When going from the loop group to the Kac-Moody group, we have two extra $\C^*$-factors- one from the central extension and one from the loop rotation. It is already in \cite{MV} that extra loop rotation equivariance doesn't change the category- the same proof shows that the same holds with the extra central $\C^*$. I can also add a formal proof if you would want me to.}

Finally, we take equivariant global sections. Let $R_{\aff}$ denote the equivariant cohomology of a point under the Kac-Moody torus $T_{\KM}$. It is a polynomial ring over a \BE{I changed ``the'' to ``a''. There are two extra copies of $\C^*$ you say? then the polynomial ring gets bigger by two variables, not one. Though the new one is central and plays no role. probably we add a remark below about this, that also includes your point about MV?} reflection representation of the affine Weyl group $W_{\aff}$, and inherits an action of $W_{\aff}$. The $P_{\fin}$-equivariant cohomology of a point becomes the subring $R_{\aff}^{W_{\fin}}$ of polynomials invariant under $W_{\fin}$. For other $\chi \in \Omega$, the $P_{\chi}$-equivariant cohomology of the point is the invariant subring $R_{\aff}^{\chi}$ under a different copy of $W_{\fin}$ inside $W_{\aff}$. Given a $(P_{\chi_1},P_{\chi_2})$-equivariant perverse sheaf on $G_{\KM}$, its equivariant global sections form a graded $(R_{\aff}^{\chi_1},R_{\aff}^{\chi_2})$-bimodule. The bimodules which arise in this fashion are known as singular Soergel bimodules, and thanks to the decomposition theorem they have an algebraic description as well (see the next section), which is what enables their study without needing to think about geometry at all.

\begin{remark} Meanwhile, the $T(\OC)$-equivariant cohomology of a point is the smaller ring $R_{\fin}$, and the $G(\OC)$-equivariant cohomology of a point is $R_{\fin}^{W_{\fin}}$. Lifting to the Kac-Moody setting provides a larger ring $R_{\aff}$ within which one can compare the equivariant cohomologies under various (lifts of) conjugates of $G(\OC)$.\end{remark}

\begin{remark} \BE{is this ok?} The decomposition theorem holds when sheaves have coefficients in a field of characteristic zero. When working in positive characteristic (but avoiding bad primes), singular Soergel bimodules correspond to parity sheaves \BE{cite JMW} rather than perverse sheaves. To prove a Soergelified Satake theorem in this setting, one would need the stronger version of our theorem discussed in Remark \ref{rmk:youthfuloptimism} above and \BE{XXXX} below. \end{remark}

% \PT{If we want to talk about sheaf coefficients in positive characteristic, we would want to work with parity sheaves. I think we also would want to avoid bad primes for each type, though in type A IC sheaves are probably parity in all characteristic outside 2. Will have to check.}



Soergel's famous Struktursatz \cite{SoerXX} states, for a simple Lie group $G$ with Borel $B$, that the equivariant global sections functor is fully faithful for \emph{semisimple} $B$-equivariant perverse sheaves on $G/B$. His student Harterich, in his thesis \cite{Harterich}, extended this to Kac-Moody groups and their flag varieties. Abe \cite{Abe} extends this further to partial flag varieties, whence the geometric side of the Satake equivalence can be replaced with singular Soergel bimodules (we omit some technical details here that will play a role in \BE{XXX} below) \BE{I'm thinking about equivalence between Abe and Williamson}.\PT{I thought this was immediate from the Hom-formula and the fact that Abe's morphisms are actually just bimodule morphisms compatible with the extra data (in particular, the forgetful functor is faithful). The Hom formula then gives us fullness because each graded Hom space is finite dimensional. Did I miss something?}
% \BE{appropriate citation or folklore}\PT{Technically, \cite{Abe} is the only place that I know where it is written up. When Williamson's categorification result is true, we have that the forgetful functor from Abe's SSBim to Williamson's SSBim is fully faithful (with faithfulness always true by definition of morphisms in Abe's category). TLDR: we now indeed do have a reference for this folklore!!} extends this further to partial flag varieties, whence the geometric side of the Satake equivalence can be replaced with singular Soergel bimodules.

\subsection{Singular Soergel bimodules}

Going forward, we restrict our attention to type $A_{n-1}$ for simplicity, where the group $\Omega$ is isomorphic to $\Z/n\Z$. Rather than using $\chi$ to represent an element of $\Omega$, we now use integers like $a, b, c$, which we identify with their residue classes modulo $n$.

Let $V$ be the $\Q$-linear reflection representation of $W_{\aff}$. Let $R = R_{\aff} := \Sym(V)$ denote the polynomial ring whose linear terms are $V$. We double the usual grading on $R$, so that linear terms live in degree $2$.

The natural action of the affine Weyl group $W_{\aff}$ on $V$ extends to an action by ring homomorphisms in $R_{\aff}$. Associated to each proper subset $I \subset S_{\aff}$ of the affine simple reflections, we have a finite parabolic subgroup $W_I \subset W_{\aff}$, and can consider the subring $R^I$ of polynomials invariant under $W_I$. Whenever $I \subset J$ we have $R^J \subset R^I$, and this ring extension is a \emph{graded Frobenius extension}, meaning that induction and coinduction (both functors from $R^J$-modules to $R^I$-modules) are equivalent functors up to a grading shift. Said another way, induction and restriction are biadjoint up to a shift. Note that induction and restriction functors are realized by tensoring with graded bimodules. Below we set
\begin{equation} \Ind^I_J := R^I \in (R^I, R^J)-\gbimod, \qquad \Res^I_J := R^I(\ell(w_J) - \ell(w_I)) \in (R^J,R^I)-\gbimod. \end{equation}
The number $\ell(w_J) - \ell(w_I)$ in parentheses for $\Res^I_J$ represents a grading shift, see \cref{subsec- SSBim} for details.

Along a sequence of proper subsets
$[[ I_0 \subset I_1 \supset I_2 \subset \ldots \supset I_d]]$
one can consider the iterated induction and (shifted) restriction functor, which corresponds to tensor product with the graded $(R^{I_0}, R^{I_d})$-bimodule
$$ R^{I_0} \otimes_{R^{I_1}} R^{I_2} \otimes \cdots \otimes R^{I_d}(k)$$
for some easily computable integer $k$. Such bimodules are called \emph{singular Bott-Samelson bimodules}, and they form a 2-category $\SBSBim$. Direct sums of direct summands of shifts of such bimodules are called \emph{singular Soergel bimodules}, and they form a Karoubian graded $2$-category $\SSBim$.  To be more precise, the objects of $\SBSBim$ (or $\SSBim$) are proper subsets $I \subset S$, and the category $\Hom(I,J)$ is a full subcategory of $(R^J,R^I)-\gbimod$. 

As noted in the previous section, there is a bijection between $\Omega$ and the isomorphic copies of the finite Dynkin diagram $S_{\fin}$ inside $S_{\aff}$. That is, for each $a \in \Omega$ we can associate a vertex in $S_{\aff}$, and the corresponding subset $\hh{a} := S_{\aff} \setminus a$ satisfies $W_{\hh{a}} \cong W_{\fin}$. The Soergelified Satake equivalence will be interested primarily in those singular Soergel bimodules which are $(R^{\hh{a}_1}, R^{\hh{a}_2})$-bimodules for $a_1, a_2 \in \Omega$.

We can deform $V$ to a $\Q(q)$-vector space $V_{q}$ with a $W_{\aff}$ action, using a deformed Cartan matrix as in \eqref{eq:deformedq}. 
(For a precise definition see \BE{ref}\BE{By the way, in comments I say ref for internal refrecne and cite for external citation.}\PT{Got it, thanks for clarifying!}) 
The entire story above can be repeated for $V_{q}$, producing a graded ring $R_q := \Sym(V_{q})$, with subrings $R^I_{q}$ for proper subsets $I \subset S_{\aff}$. 
One obtain 2-categories $\SBSBim_{q}$ and $\SSBim_{q}$ of graded bimodules over various invariant subrings.

Within the monoidal category \BE{I removed subscripts $\Bbbk$ on the reps here} $\Rep \slf_n$ one can find a (non-additive) monoidal subcategory $\Fund \slf_n$, whose objects are iterated tensor products of fundamental representations\footnote{A fundamental representation is an irreducible representations whose highest weight is a fundamental weight.}. Just as the category $\SSBim$ is obtained as the Karoubi envelope of the strict 2-category $\SBSBim$, the 2-category $\Rep^\Omega = \Rep^{\Omega} \gf^{\vee}$ can be viewed as the Karoubi envelope of the strict 2-category $\Fund^{\Omega}$. These categories also have quantum group analogues, which we denote $\Rep^\Omega_{q}$ and  $\Fund_{q}^{\Omega}$.

The Soergelified Satake equivalence is the $2$-functor in the following theorem.

\begin{theorem} \label{thm:main} Let $\Fund_q^{\Omega} := \Fund^{\Omega} U_q(\slf_n)$. There is a strict 2-functor $\GS_{q} \colon \Fund^{\Omega}_q \to \SBSBim_{q}$. On objects it sends $a \mapsto \hh{a}$. The $2$-category $\Fund_{q}^{\Omega}$ is generated by fundamental representations $L_b$ of highest weight $\varpi_b$, viewed as $1$-morphisms from $\Rep_a$ to $\Rep_{a+b}$, for various $a, b \in \Omega$ with $b \ne 0$. This $1$-morphism is sent to the graded $(R^{\hh{a+b}},R^{\hh{a}})$ bimodule \BE{get shift right}
\[ L_b \mapsto R^{\hh{a}, \hh{a+b}}(k).\]
That is, $L_b$ is sent to the singular Bott-Samelson bimodule associated to the sequence 
\[ [[S_{\aff} \setminus \{a+b\} \supset S_{\aff} \setminus \{a,a+b\} \subset S_{\aff} \setminus \{a\}]],\] the composition of induction followed by restriction.
% a composition of induction from $R^{\hh{a}}$ to $R^{\hh{a}, \hh{a+b}}$, followed by (shifted) restriction from $R^{\hh{a}, \hh{a+b}}$ to $R^{\hh{a+b}}$.
What $\GS_q$ does to $2$-morphisms will be explained in \BE{ref}. \end{theorem}
\PT{$\GS_q$? I can add a discussion showing that our theorem for $\GS_\zeta$ has a $\Q(q)$ form, if we want the $q-$version as well. For this the caps and cups would be scaled according to the difference in the Frobenius structure.}\BE{yes, please}

\begin{remark} Outside of type $A$, fundamental representations are not typically sent to Bott-Samelson bimodules, but to indecomposable Soergel bimodules which are summands thereof. \end{remark}

\subsection{Diagrammatics}

\BE{Note to self - stop procrastinating by reading the introduction, go to the body. But when you return to intro, return here.}

Let $\Bbbk$ denote the base field over which we take our representations. In type $A$ fundamental representations are special in that they are tilting in any characteristic (as are their tensor products). A consequence of this is that $\Fund \slf_n$ is flat: the size of morphism spaces between tilting modules does not depend on the characteristic of the ground field, and is unaffected by base change. One might then expect an integral form $\Fund_{\Z} \slf_n$ to exist: a $\Z$-linear monoidal category which becomes isomorphic to $\Fund_{\Bbbk} \slf_n$ after base change to any field $\Bbbk$.

The standard way to produce an integral form of an algebra or category is to provide a presentation by generators and relations, where all the coefficients in the relations are integers. Such a presentation was provided as the category of $\slf_n$ webs, due to \cite{Kupe, CKM}. This presentation has been refined in various iterations, and we use the version found in \cite{Bigelow} \BE{cite your fav source philip} \PT{Yup.}. It was proven in \cite{ELLCC} that webs do indeed provide a suitable integral form for $\Fund \slf_n$. Modifying these statements to account for the 2-category $\Fund^{\Omega} \slf_n \subset \Rep^{\Omega} \slf_n$ is straightforward, involving the use of \emph{colored webs}, see \cite{EQuantumI} (or below).

On the other side of the equivalence, it is known that singular Soergel bimodules are not flat: in finite characteristic the reflection representation of the affine Weyl group is no longer faithful, and there are additional morphisms between singular Soergel bimodules. One seeks a presentation by generators and relations of $\SBSBim$, and one has not yet been discovered. There is a lot of progress towards this goal, however. Let us discuss the state of the art, as of the writing of this paper.

In \cite{EWSFrob} one can find a general diagrammatic scheme for morphisms between tensor products of induction and restriction bimodules, when one is given a collection of Frobenius extensions indexed by a cube. In this case, the cube in question is the set of subsets of $S_{\aff}$; one must only consider proper subsets, which does not cause any issues. The diagrammatic calculus of \cite{EWSFrob} has enough generating 2-morphisms, and has a number of relations which hold true for any cube of Frobenius extensions, but is known to be missing relations which are special to e.g. the setting of singular Soergel bimodules associated to a Coxeter system. Let us call this diagrammatic category $\DiagSBS'$; it has a functor $F'$ to singular Soergel bimodules. Note that, like any category defined by generators and relations, $\DiagSBS'$ has an integral form defined over a small extension of the integers (containing the elements of the Cartan matrix).

There should be a diagrammatic category $\DiagSBS$ which is a quotient of $\DiagSBS'$ (with additional, unknown relations imposed), so that the functor $F'$ factors through this quotient. The induced functor from $\DiagSBS$ to $\SBSBim$ is denoted $F$, and it should be the case that $F$ is an equivalence after base change (under certain assumptions on the base ring). Moreover, $\DiagSBS$ should be flat, with morphism spaces whose sizes are governed by the Soergel-Williamson Hom formula (see \BE{ref}\PT{\cref{thm-hom formula}?}). Such a 2-category $\DiagSBS$ is sure to exist, so we speak as though it exists, but the precise relations one needs to impose are unknown.

In \cite{EKLP}, a collection of diagrams in $\DiagSBS'$ are found which, after applying the functor $F'$, are sent to a basis for morphisms between singular Bott-Samelson bimodules. These diagrams are called \emph{double leaves}. Consequently, double leaves will descend to a basis for $\DiagSBS$, once it is properly defined.

The following is an upgrade of Theorem \ref{thm:main}.

\begin{conjecture} \label{conj:diag} There is a strict 2-functor $\GS_q \colon \Fund_{q}^{\Omega} \to \DiagSBS_{q}(\tilde{A}_{n-1})$, defined over the base ring $\Z[q,q^{-1}]$, through which the 2-functor of Theorem \ref{thm:main} factors. It is a degree zero equivalence. \end{conjecture}

Conjecture \ref{conj:diag} would let one algebraically explore the geometric Satake equivalence in finite characteristic, or its quantum analogue when $q$ is a root of unity.

What this functor $\GS_q$ does to diagrams was explicitly stated in \cite{EQuantumI} using the language of $\DiagSBS'$ \PT{Precise reference?}, though it is clearly not a functor to $\DiagSBS'$; one needs the additional relations of $\DiagSBS$ in order for the relations between webs to be sent to zero. To prove Conjecture \ref{conj:diag} the remaining questions are: once additional relations are imposed, can one check the relations between webs, so that the 2-functor is well-defined? Will it be fully faithful to degree zero?

Note that ordinary geometric Satake implies that certain pairings in the affine Hecke algebra match the dimensions of morphism spaces in $\Fund^{\Omega}$, and by the Soergel-Williamson hom formula this implies that the dimensions of morphism spaces (in degree zero) agree on both sides of the functor $\GS_q$. 
The question of whether $\GS_q$ is fully faithful to degree zero reduces to the question of fullness. \BE{I completely forget if we've discussed this yet or what your plans were. One can prove this by showing that light leaves are the images of light ladders, and this is a nice thing to do anyway... if we want to add that to the paper.}\PT{Probably not in this paper, since it's already pretty big.}


In \cite{ECathedral}\PT{Updated reference to Dihedral Cathedral}, the additional relations were found for dihedral groups, producing the category $\DiagSBS$. When $n= 2$ the affine Weyl group is dihedral. More generally, there are only expected to be relations associated to finite parabolic subgroups of the Coxeter group in question. When $n=3$ all finite subgroups of the affine Weyl group are dihedral, and it was proven in \cite[Appendix]{EQuantumI} that the expected relations did in fact suffice to produce $\DiagSBS$ for $\tilde{A}_{2}$. Using this, \cite{EQuantumI} was above to prove Conjecture \ref{conj:diag} for $n=2, 3$. 

\begin{remark} Diagrammatics for $n=4$ are known to the experts but not yet in the literature, and Elias was also able to prove Conjecture \ref{conj:diag} for $n=4$, in unpublished computations. \end{remark}

Beyond small rank, there has been a great deal of recent progress in the study of $\DiagSBS$. In type $A$, one expects that the only additional relations one needs to present $\DiagSBS$ are the categorified MOY relations for finite symmetric groups. On the other side of Schur-Weyl duality, these relations are categorified by the Stosic formula found in \cite{KLMS}. The Stosic formula can be transferred to $\DiagSBS'$ using a 2-functor found in \cite{ELosev}. Proving that double leaves is indeed a basis after imposing these relations is an active research program of the authors of \cite{EKLP}.

When \cite{EQuantumI} first appeared, it was erroneously assumed that $\DiagSBS$ was not long to be awaited, and that once it was invented, some quick computations would prove Conjecture \ref{conj:diag} in general. Sadly, this was not the case, nor are  the computations quite so easy.

As a consequence, we abandon the goal of proving Conjecture \ref{conj:diag}, and focus our attention on proving Theorem \ref{thm:main}. We have a potential 2-functor $\GS_q$ (defined via diagrams in $\DiagSBS'$ and the map $F'$) and need only check the relations and prove it is a degree zero equivalence. This should be easier, because singular Soergel bimodules already are known to satisfy the Soergel-Williamson hom formula thanks to \BE{cite}.

Actually, that last statement is true when $q=1$, but is not yet in the literature in general! When $q \ne 1$ the Cartan matrix \eqref{eq:deformedq} is \emph{unbalanced}, meaning that roots are only well-defined up to scalar. Williamson's results from \cite{WillSingular} tacitly assume a balanced realization. Thankfully, we have combed through his work, dotted the $i$s and crossed the $t$s, and verified that most of it goes through verbatim; we provide the necessary modifications in \cref{subsec-categorification} \BE{ref}\PT{Added.}.




, \BE{had to stop for the day. next up: how one can prove relations in bimodules without needing to use diagrams, but only when the soergel hom formula holds, and with a few calculations. Then: the subtleties of unbalanced bullshit.}


\PT{Me starting to write up the rest of the introduction from here.}
Since a purely diagrammatic proof is not yet possible due to the reasons discussed before, to prove \cref{thm:main}, we explicitly check the required relations in the corresponding algebraic categories of honest bimodules.
Doing this requires various explicit computations involving invariant subrings of $R$, and this is the main content of \cref{subsec-further frob calculations}.












\PT{Still working on this. This will be the end of my addition to intro.}

Before the discuss the results in more detail, we note a change of normalization. We use a parameter $\zeta$ for which $\zeta^n=q^2$, and we use a $\zeta$-deformed Cartan matrix (below for $\tilde{A}_3$).
\begin{equation} \label{eq:deformedzeta} \begin{bmatrix}
    2 &-\zeta &0 & -\zeta^{-1} \\
    -\zeta^{-1} &2 &-\zeta &0\\
    0 &-\zeta^{-1} &2 &-\zeta \\
    -\zeta & 0&-\zeta^{-1} &2
\end{bmatrix} \end{equation}
This Cartan matrix can be obtained from \eqref{eq:deformedq} by rescaling the roots and coroots, so the corresponding categories of singular Soergel bimodules are equivalent. The $\zeta$-deformed Cartan matrix admits additional symmetries which make it easier to check the web relations.






Same as in \cite{EQuantumI}, but with a different realization for the affine Weyl group. We consider a realization with the $\zeta$-deformed Cartan matrix:

$$\begin{bmatrix}
    2 &-\zeta &0 &  &\dots &  &-\zeta^{-1} \\
    -\zeta^{-1} &2 &-\zeta &0 & &\dots &0\\
    0 &-\zeta^{-1} &2 &-\zeta &0 &\ldots &0\\
    \vdots & &\ddots &\ddots &\ddots & &\vdots\\
    \vdots & &   &-\zeta^{-1} &2 &-\zeta &0\\
    \vdots & &   & &-\zeta^{-1} &2 &-\zeta \\
    -\zeta & &\dots & & &-\zeta^{-1} &2 
\end{bmatrix}$$
where $\zeta^n=q^2$. The main result of this paper is generalizing the construction from \cite{EQuantumI} to type $A_n$ for $n\geq 4$.
\PT{@Ben I haven't touched this section since you wanted to do it, iirc.}


%%% end comment
}